% AUTORIDAD DE FUENTE: Exposé IX (12.8.9)--12.10, líneas francesas
% 3796--3950; índices de escaneo 491--495, folios 480--484.

\begin{equation}\tag{12.8.9}
\begin{tikzcd}[row sep=small,column sep=small]
0 \arrow[r] & P \arrow[rr] \arrow[d,dash,"\wr"]
 & & F \arrow[r] \arrow[d,dash,"\wr"] & R \arrow[r] & 0 & \\
& T_\ell(P)^t\simeq
  \underline{\operatorname{Hom}}
    \bigl(M,T_\ell(\sisi{\underline{\mathbb G}}{\mathbb G}_{m\underline S})\sisi{}{\bigr)}
& & T_\ell(P)^f & & & \\
& & & & & & \\
0 \arrow[r] & R \arrow[rr] \arrow[d,dash,"\wr"]
& & E \arrow[r] & Q \arrow[r] \arrow[d,dash,"\wr"] & 0 & \\
& D(R) & & & M\otimes\mathbb Z_{\ell\,\underline S} & & .
\end{tikzcd}
\end{equation}

\noindent
Obsérvese que $R$ es canónicamente autodual, como se sigue de la
autodualidad (12.8.7) y de la afirmación de ortogonalidad señalada antes.
La única diferencia con la construcción citada es que el emparejamiento
obtenido sobre $M_\ell\times M_\ell$ toma sus valores en el grupo
\[
 \bigl(\mathbb G_m(U)/\mathbb G_m(\underline S)\bigr)
 \otimes_{\mathbb Z}\mathbb Z_\ell,
\]
que ya no es canónicamente isomorfo a $\mathbb Z_\ell$, sino a
$\mathbb Z_\ell^{\,L}$, mediante el isomorfismo

\begin{equation}\tag{12.8.10}
 \mathbb Z^L\simeq
 \mathbb G_m(U)/\mathbb G_m(\underline S).
\end{equation}

\noindent
Este envía el elemento indexado por $\lambda$ de la base canónica de
$\mathbb Z^L$ a la clase de $T_{i_\lambda}$ en el segundo miembro, donde
$T_{i_\lambda}$ es una ecuación del divisor $D_\lambda$. Obtenemos así
un emparejamiento canónico

\begin{equation}\tag{12.8.11}
 \underline{u}_\ell:M_\ell\otimes M_\ell
 \longrightarrow\mathbb Z_\ell^{\,L}\qquad .
\end{equation}

\noindent
Es esencialmente inmediato que el emparejamiento de monodromía $u_\ell$
(12.4.7) sobre $S$ que nos proponemos estudiar es el compuesto de
$\underline{u}_\ell$ con el homomorfismo
$\mathbb Z_\ell^{\,L}\longrightarrow\mathbb Z_\ell$ obtenido, por
producto tensorial con $\mathbb Z_\ell$ sobre $\mathbb Z$, a partir del
homomorfismo

\begin{center}
\begin{tikzcd}
 \mathbb G_m(U)/\mathbb G_m(\underline S)
   \arrow[r] \arrow[d,dash,"\wr"'] &
 \mathbb G_m(\eta)/\mathbb G_m(S)
   \arrow[d,dash,"\wr"'] \\
 \mathbb Z^L & \mathbb Z
\end{tikzcd}
\end{center}

\noindent
inducido por $g:S\longrightarrow\underline S$ (que envía $\eta$ en
$U$). Este homomorfismo está dado por un elemento

%%%%%%%%%%  scan idx 492  |  folio 481  %%%%%%%%%%

\begin{equation}\tag{12.8.12}
 (d_\lambda)_{\lambda\in L}\in\mathbb Z^L,\qquad
 d_\lambda\geqslant1\quad\text{para todo }\lambda\in L,
\end{equation}

\noindent
donde, para cada $\lambda\in L$, $d_\lambda$ es la multiplicidad de la
imagen inversa en $S$ del divisor $D_\lambda$ de $\underline S$. Depende,
de hecho, únicamente del completado
$\widehat{\underline{\mathcal O}}_{X,x_\lambda}$ como $V$-álgebra
(Exposé VI\sisi{\hspace{3em}}{}). Además, se sabe (Exposé
VI\sisi{\hspace{3em}}{}) que, para un $\lambda$ dado, se tiene
$d_\lambda=1$ si y solo si $\sisi{x}{X}$ es \underline{regular} en
$x_\lambda$. Como en 12.5 hemos supuesto que $X$ es regular, es decir,
que todos los $d_\lambda$ son iguales a $1$, la aserción 12.5 se sigue
del lema siguiente, en el que ya no intervienen $X,S,g$.

\medskip
\noindent
\underline{Lema} 12.8.13. \underline{Si $\lambda\in L$, el compuesto de
(12.8.11) con $\operatorname{pr}_\lambda$ es igual a la forma
$\varphi_{x_\lambda}\otimes\varphi_{x_\lambda}$ de (12.4.4).}

\medskip
Para demostrar esta última aserción, consideremos el localizado estricto
completado $\underline S^\lambda$ de $\underline S$ en el punto genérico
de $\sisi{D}{D_\lambda}$, y pongamos
\[
 \underline X^\lambda=
 \sisi{X}{\underline X}\times_{\sisi{S}{\underline S}}\underline S^\lambda.
\]
La fibra especial $\underline X_o^\lambda$ tiene entonces un único punto
singular $y_\lambda$, que sigue siendo una singularidad cuadrática
ordinaria. En el emparejamiento de monodromía correspondiente

\begin{equation}\tag{12.8.14}
 u_\ell^\lambda:M_\ell^\lambda\otimes M_\ell^\lambda
 \longrightarrow\mathbb Z_\ell\qquad ,
\end{equation}

\noindent
$M^\lambda$ tiene, por consiguiente, rango $1$ o $0$. Dejamos al lector
la verificación de los hechos sencillos siguientes, que explicitan, en
el caso particular en que los necesitamos, el comportamiento de las
construcciones precedentes bajo localización.

%%%%%%%%%%  scan idx 493  |  folio 482  %%%%%%%%%%

\medskip
\noindent
a) Existe un único subtoro $T_\lambda$ de
$P\times_{\underline S}D_\lambda$ cuya restricción al punto genérico
$s_\lambda$ de $D_\lambda$ es el toro maximal de $P_{s_\lambda}$ (que
tiene dimensión $0$ o $1$).

\medskip
\noindent
b) Consideremos el homomorfismo transpuesto $M\longrightarrow M^\lambda$
sobre los grupos de caracteres de la inclusión
$T_{\lambda o}\hookrightarrow T_o$, donde $T_o$ es el toro maximal de
$P_o$, y la forma
$\varphi_{y_\lambda}:M^\lambda\longrightarrow\mathbb Z$, definida salvo
signo como en 12.4.3 (sustituyendo $X$ por $\underline X^\lambda$).
Entonces, salvo signo, el diagrama siguiente es conmutativo:

\begin{equation}\tag{12.8.15}
\begin{tikzcd}[column sep=huge,row sep=large]
 M \arrow[rr] \arrow[dr,"\varphi_{x_\lambda}"']
 & & M^\lambda \arrow[dl,"\sisi{\varphi_{\lambda_\lambda}}{\varphi_{y_\lambda}}"] \\
 & \mathbb Z & \qquad .
\end{tikzcd}
\end{equation}

\noindent
(En efecto, elegir un orden de las dos ramas de $X_o$ que pasan por
$x_\lambda$ determina una elección análoga en el punto $y_\lambda$ de
$\underline X_o^\lambda$ y, para las elecciones así correspondientes,
el diagrama precedente es conmutativo.)

%%%%%%%%%%  scan idx 494  |  folio 483  %%%%%%%%%%

\medskip
\noindent
c) El diagrama siguiente es conmutativo:

\begin{center}
\begin{tikzcd}[column sep=huge,row sep=large]
 M_\ell\otimes M_\ell
   \arrow[rr] \arrow[dr,"\operatorname{pr}_\lambda\underline{u}_\ell"']
 & & M_\ell^\lambda\otimes M_\ell^\lambda
   \arrow[dl,"u_\ell^\lambda"] \\
 & \mathbb Z_\ell & \qquad ,
\end{tikzcd}
\end{center}

\noindent
donde la notación es la de 12.8.13 y (12.8.14). Estas relaciones muestran
inmediatamente que, para demostrar 12.8.13, basta demostrar la relación
análoga para $\underline X^\lambda$:
\[
 u_\ell^\lambda=
 (\varphi_{y_\lambda}\otimes\varphi_{y_\lambda})
 \otimes_{\mathbb Z}\mathbb Z_\ell,
\]
que no es otra cosa que 12.5 para $\underline X^\lambda$ sobre
$\underline S^\lambda$. Ahora bien, $\sisi{\underline S}{\underline
S^\lambda}$ es un trait de \underline{característica cero}. Esto concluye
la reducción anunciada.

\medskip
\noindent
\underline{Observaciones} 12.9. Obsérvese que esta reducción nos reduce
simultáneamente al caso en que $X_o$ tiene un solo punto singular, que es
el caso considerado clásicamente en la fórmula de Picard--Lefschetz.
Obsérvese también que, en este caso, $X_o$ tiene una o dos componentes
irreducibles, correspondientes respectivamente a
$\operatorname{rank}M=1$ y $\operatorname{rank}M=0$ (este último es el
caso en que la jacobiana de la fibra genérica tiene buena reducción).

%%%%%%%%%%  scan idx 495  |  folio 484  %%%%%%%%%%

\medskip
\noindent
12.10. La demostración dada en 12.8 proporciona un resultado más general
que 12.5, válido sin suponer que $X$ sea regular, introduciendo los
enteros $d_\lambda\geqslant1$ de (12.8.12): todos los emparejamientos de
monodromía se deducen entonces del emparejamiento

\begin{equation}\tag{12.10.1}
 u=\sum_{\lambda\in L}d_\lambda\,
   \varphi_{x_\lambda}\otimes\varphi_{x_\lambda}:
   M\otimes M\longrightarrow\mathbb Z\qquad .
\end{equation}
